% 前置：题献 + 序言（Preface）
\chapter*{序言}
\addcontentsline{toc}{chapter}{序言}

\begin{center}
{\kaishu 致我的父母 Ruth 与 Israel}
\end{center}

\medskip

在紧张而迅速的进展之中，为大多数研究者而言，撰写教学用书的优先级并不高。然而，在把我的讲义\footnote{为 1990--1993 年间在波士顿大学与以色列理工学院开设的研究生课程“量子多体系统”而写的讲义。}改写成这本书的过程中，我发现了一个个人的收益：把我所理解的内容按一种（但愿是）简单的逻辑次序组织起来。本书中属于我本人原创贡献的内容很少，大部分知识取自研究文献，还有一些来自与同事们的交谈——一种在同门同道之间口耳相传的物理学口述传统。

多年来，图微扰论一直是处理金属、半导体、巡游磁体与超导体中相互作用的主要理论工具。它本质上是一种围绕自由准粒子的弱耦合展开。过去十年间的许多实验发现——包括重费米子、分数量子霍尔效应、高温超导与量子自旋链——从弱耦合的观点并不容易理解。因此，近年来，用于处理强电子--电子相互作用的另一类非微扰工具得到了蓬勃的发展。

本书集中于强关联（或受约束）量子系统的两个基本范式：Hubbard 模型与 Heisenberg 模型。这两个模型承载了许多基本概念，例如有效哈密顿量、变分基态、自发对称性破缺与量子无序；此外，它们也被用作各种非微扰近似方案的试验场，这些方案在理论物理的众多领域中都有应用。

本书的程度适合具备一定固体物理背景（单电子理论）并熟悉二次量子化的研究生。习题难度不一，以具体例子和推论补充正文。部分数学背景材料放在附录中。

我最要感谢的是 Maxim Raykin、Moshe Havilio 与 Ziad Musslimani 孜孜不倦的努力，他们细致的校对清除了一处处不一致，并帮助澄清了众多论点。我还深深受益于 Duncan Haldane，是他把我领进量子磁学的大门；也感谢我的朋友和同事 Dan Arovas，他教会我母哈密顿量、单模近似以及许多其他东西\footnote{包括幻匕首（phantom daggers）的用法。}，并提出了许多批评意见。感谢 Alfred P. Sloan 基金会的支持，使我得以完成本书。

\begin{flushright}
Assa Auerbach\\
海法，1997
\end{flushright}
